\documentclass[11pt,a4paper]{article}
\usepackage[T1]{fontenc}\usepackage[utf8]{inputenc}\usepackage[czech]{babel}
\usepackage{lmodern,microtype,amsmath,amssymb,geometry,enumitem}
\geometry{top=17mm,bottom=17mm,left=22mm,right=22mm}
\pagestyle{empty}\setlength{\parindent}{0pt}\setlength{\parskip}{3pt}
\begin{document}
\begin{center}\Large\bfseries 95. hodina -- řešení sbírky\end{center}
\small Řešení jsou určena učiteli. U úloh s parametrem je uvedena úplná diskuse; u ostatních zadání ověřte i původní podmínky.\normalsize\par\medskip
\noindent\textbf{Úloha 1 (A).} Vyřešte $4(x-2)=2x+6$.
\[\begin{aligned}
4x-8&=2x+6 \\
2x&=14 \\
x&=7 \\
K&=\{7\}
\end{aligned}\]
\noindent\textbf{Úloha 2 (A).} Vyřešte $\frac{2x-1}{3}-\frac{x+2}{2}\le1$.
\[\begin{aligned}
2(2x-1)-3(x+2)&\le6 \\
x-8&\le6 \\
x&\le14 \\
K&=(-\infty;14\rangle
\end{aligned}\]
\noindent\textbf{Úloha 3 (A).} Vyřešte $\begin{cases}x+y=7,\\2x-y=2.\end{cases}$
\[\begin{aligned}
3x&=9 \\
x&=3,\quad y=4 \\
K&=\{[3;4]\}
\end{aligned}\]
\noindent\textbf{Úloha 4 (A).} Vyřešte $|x-3|=4$.
\[\begin{aligned}
x-3&=\pm4 \\
x&=-1\ \text{nebo}\ x=7 \\
K&=\{-1;7\}
\end{aligned}\]
\noindent\textbf{Úloha 5 (A).} Vyřešte $|x-2|\le3$.
\[\begin{aligned}
-3&\le x-2\le3 \\
-1&\le x\le5 \\
K&=\langle-1;5\rangle
\end{aligned}\]
\noindent\textbf{Úloha 6 (B).} Vyřešte $3[2x-(x-4)]=2x+16$.
\[\begin{aligned}
3(x+4)&=2x+16 \\
3x+12&=2x+16 \\
x&=4 \\
K&=\{4\}
\end{aligned}\]
\noindent\textbf{Úloha 7 (B).} Vyřešte $2-3(x-1)>x+5$.
\[\begin{aligned}
5-3x&>x+5 \\
-4x&>0 \\
x&<0 \\
K&=(-\infty;0)
\end{aligned}\]
\noindent\textbf{Úloha 8 (B).} Vyřešte $\begin{cases}2x-3y=7,\\4x+y=9.\end{cases}$
\[\begin{aligned}
y&=9-4x \\
2x-3(9-4x)&=7 \\
14x&=34 \\
x&=\frac{17}{7},\quad y=-\frac57 \\
K&=\left\{\left[\frac{17}{7};-\frac57\right]\right\}
\end{aligned}\]
\noindent\textbf{Úloha 9 (B).} Vyřešte $|x-2|+|x+2|=6$.
\[\begin{aligned}
x<-2&:\quad -2x=6\Rightarrow x=-3 \\
-2\le x\le2&:\quad4\ne6 \\
x>2&:\quad2x=6\Rightarrow x=3 \\
K&=\{-3;3\}
\end{aligned}\]
\noindent\textbf{Úloha 10 (B).} Diskutujte $(a-2)x=a+4$ pro všechna $a\in\mathbb R$.
\[\begin{aligned}
a\ne2&:\quad x=\frac{a+4}{a-2} \\
a=2&:\quad0=6\ \text{-- spor} \\
K_a&=\begin{cases}\left\{\frac{a+4}{a-2}\right\},&a\ne2,\\\varnothing,&a=2.\end{cases}
\end{aligned}\]
\noindent\textbf{Úloha 11 (B).} Ze vzorce $Q=cm(T_2-T_1)$ vyjádřete $T_2$ včetně podmínky.
\[\begin{aligned}
Q&=cm(T_2-T_1) \\
T_2&=T_1+\frac{Q}{cm}\quad(cm\ne0)
\end{aligned}\]
\noindent\textbf{Úloha 12 (C).} Vyřešte $\begin{cases}x+y+z=6,\\2x-y+z=3,\\x+2y-z=2.\end{cases}$
\[\begin{aligned}
(\mathrm{II})-(\mathrm I)&:\quad x-2y=-3 \\
(\mathrm I)+(\mathrm{III})&:\quad2x+3y=8 \\
x&=2y-3 \\
7y&=14 \\
x&=1,\ y=2,\ z=3 \\
K&=\{[1;2;3]\}
\end{aligned}\]
\noindent\textbf{Úloha 13 (C).} Diskutujte $(a-1)x=a^2-1$ pro všechna $a\in\mathbb R$.
\[\begin{aligned}
(a-1)x&=(a-1)(a+1) \\
a\ne1&:\quad x=a+1 \\
a=1&:\quad0=0 \\
K_a&=\begin{cases}\{a+1\},&a\ne1,\\\mathbb R,&a=1.\end{cases}
\end{aligned}\]
\noindent\textbf{Úloha 14 (C).} Vyřešte $\frac{x-1}{x+2}\ge0$.
\[\begin{aligned}
x&\ne-2 \\
x-1&=0\Rightarrow x=1 \\
\text{znaménka }&+,-,+\quad\text{na }(-\infty;-2),(-2;1),(1;+\infty) \\
K&=(-\infty;-2)\cup\langle1;+\infty)
\end{aligned}\]
\noindent\textbf{Úloha 15 (C).} Tarif A stojí $100+4m$ Kč, B stojí $220+2m$ Kč. Při kolika minutách stojí stejně? Který je levnější pro menší $m$?
\[\begin{aligned}
100+4m&=220+2m \\
2m&=120 \\
m&=60\,\mathrm{min}
\end{aligned}\]
\noindent Pro $0\le m<60$ je levnější A, při $m>60$ je levnější B; v $m=60$ stojí oba 340 Kč.
\noindent\textbf{Úloha 16 (C).} Vyřešte $|x-2|+|x+2|=4$.
\[\begin{aligned}
x<-2&:\quad-2x=4\Rightarrow x=-2\not<-2 \\
-2\le x\le2&:\quad4=4 \\
x>2&:\quad2x=4\Rightarrow x=2\not>2 \\
K&=\langle-2;2\rangle
\end{aligned}\]
\end{document}
